\chapter{Group Theory}
\label{ch:groups}

\medskip

\textbf{Notation.} $G$ --- group; $H\le G$ --- subgroup; $H\trianglelefteq G$ --- normal subgroup.
$[G:H]$ --- index; $|G|$ --- order.
$Z(G)$ --- center; $C_G(H)$ --- centralizer; $N_G(H)$ --- normalizer.
$\operatorname{cl}(g)$ --- conjugacy class of $g$.
$\operatorname{Aut}(G)$ --- automorphism group; $\operatorname{Inn}(G)$ --- inner automorphisms.
$G/H$ --- quotient group; $G\cong H$ --- isomorphic.
$S_n$ --- symmetric group; $A_n$ --- alternating group.
$\mathbb Z/n\mathbb Z$ (or $C_n$) --- cyclic group of order $n$.

\medskip

\section{Orders, Lagrange, Cauchy, cyclic groups}
\label{sec:groups-orders}

\subsection{Lagrange's theorem}
\label{subsec:lagrange}

$|G|<\infty$, $H\le G$ $\Rightarrow$ $|H|\mid|G|$ and $|G|=[G:H]\cdot|H|$. Corollary: $\operatorname{ord}(g)\mid|G|$, $g^{|G|}=e$.

\subsection{Cauchy's theorem}
\label{subsec:cauchy-group}

$p$ prime, $p\mid|G|$ $\Rightarrow$ $\exists$ element of order $p$.

\subsection{Order of a power}
\label{subsec:order-power}

$\operatorname{ord}(g)=n<\infty$ $\Rightarrow$ $\operatorname{ord}(g^k)=n/\gcd(k,n)$.

\subsection{Subgroup structure of cyclic groups}
\label{subsec:cyclic-subgroups}

$G=\langle g\rangle$, $|G|=n$: every subgroup is cyclic; for each $d\mid n$ there is exactly one subgroup of order $d$: $\langle g^{n/d}\rangle$.

\subsection{Order of product in abelian groups}
\label{subsec:order-product-abelian}

$G$ abelian, $\operatorname{ord}(a)=m$, $\operatorname{ord}(b)=n$:
\begin{enumerate}
\item $\gcd(m,n)=1$ $\Rightarrow$ $\operatorname{ord}(ab)=mn$.
\item In $\langle a,b\rangle$ $\exists$ an element of order $\operatorname{lcm}(m,n)$.
\end{enumerate}
Caution: in non-abelian groups this may fail ($S_3$).

\subsection{Cyclicity criterion}
\label{subsec:cyclicity-criterion}

A finite $G$ is cyclic $\Leftrightarrow$ for each $d\mid|G|$ the equation $x^d=e$ has at most $d$ solutions. Equivalently: for each $d$ at most one subgroup of order $d$.

\subsection{Groups of order $p^2$ are abelian}
\label{subsec:order-p2}

$|G|=p^2$ $\Rightarrow$ $G$ is abelian. Up to isomorphism: $\mathbb Z/p^2$ and $\mathbb Z/p\times\mathbb Z/p$.

\subsection{Groups of order $pq$}
\label{subsec:order-pq}

$|G|=pq$, $p<q$ primes. Then:
\begin{enumerate}
\item $p\nmid q-1$ $\Rightarrow$ $G\cong\mathbb Z/(pq)$.
\item $p\mid q-1$ $\Rightarrow$ two groups: cyclic and $\mathbb Z/q\rtimes\mathbb Z/p$ (non-abelian semidirect product $G=N\rtimes H$ where $N\trianglelefteq G$, $NH=G$, $N\cap H=\{e\}$).
\end{enumerate}

\subsection{Subgroup of smallest prime index is normal}
\label{subsec:smallest-prime-index}

$[G:H]=p$ is the smallest prime divisor of $|G|$ $\Rightarrow$ $H\trianglelefteq G$.

\subsection{Formula $|HK|$}
\label{subsec:hk-formula}

$H,K\le G$ finite $\Rightarrow$ $|HK|=|H|\cdot|K|/|H\cap K|$.

\subsection{Index of intersection}
\label{subsec:index-intersection}

$[G:H\cap K]\le[G:H]\cdot[G:K]$.

\subsection{Chinese remainder for cyclic groups}
\label{subsec:crt-cyclic}

$\gcd(m,n)=1$ $\Rightarrow$ $\mathbb Z/m\times\mathbb Z/n\cong\mathbb Z/(mn)$.

\subsection{$g^2=e$ everywhere $\Rightarrow$ $G$ is abelian}
\label{subsec:g2-e}

If $\forall g\in G\;g^2=e$, then $G$ is abelian. If $G$ is finite --- $G\cong(\mathbb Z/2)^k$.

\medskip

\section{Group action: orbit--stabilizer, class equation}
\label{sec:group-action}

\subsection{Orbit--stabilizer theorem}
\label{subsec:orbit-stabilizer}

$G$ acts on $X$, $x\in X$. The map $gG_x\mapsto gx$ is a bijection $G/G_x\to Gx$. When $|G|<\infty$:
$$|Gx|=[G:G_x]=\frac{|G|}{|G_x|}.$$

\subsection{Class equation (conjugation)}
\label{subsec:class-equation}

$G$ acts on itself by conjugation $g\cdot h=ghg^{-1}$.
\begin{enumerate}
\item $|[h]|=[G:C_G(h)]$.
\item $|G|=|Z(G)|+\sum_i[G:C_G(g_i)]$ (over non-central classes).
\item The number of conjugates of a subgroup $H$ equals $[G:N_G(H)]$.
\end{enumerate}

\subsection{Fixed point lemma for $p$-groups}
\label{subsec:pgroup-fixed}

$P$ is a $p$-group acting on a finite $X$ $\Rightarrow$ $|X^P|\equiv|X|\pmod p$. If $p\nmid|X|$, there is a fixed point.

\subsection{Center of a non-trivial finite $p$-group is non-trivial}
\label{subsec:pgroup-center}

$|G|=p^n$, $n\ge1$ $\Rightarrow$ $p\mid|Z(G)|$ (in particular $Z(G)\neq\{e\}$).

\subsection{Burnside's lemma (Cauchy--Frobenius)}
\label{subsec:burnside}

$$|X/G|=\frac1{|G|}\sum_{g\in G}|X^g|,$$
where $X^g$ are the fixed points of $g$. P\'olya variant: $G\le S_n$, number of orbits on functions $[n]\to Y$, $|Y|=k$: $\frac1{|G|}\sum_g k^{c(g)}$, $c(g)$ is the number of cycles of $g$.

\subsection{Cayley's theorem (generalized)}
\label{subsec:cayley}

\begin{enumerate}
\item $G$ embeds into $S_{|G|}$ (by left translation).
\item If $[G:H]=n$, the action on $G/H$ gives $\varphi:G\to S_n$ with $\ker\varphi=\bigcap_{g}gHg^{-1}$. The kernel is a normal subgroup inside $H$, $[G:\ker\varphi]\mid n!$.
\end{enumerate}

\medskip

\section{Conjugacy classes, normalizer, centralizer}
\label{sec:conjugacy}

\subsection{Class size divides index of center}
\label{subsec:class-size-center}

$|\operatorname{cl}(g)|\mid[G:Z(G)]$. $|\operatorname{cl}(g)|=1 \iff g\in Z(G)$.

\subsection{$G/Z(G)$ cyclic $\Rightarrow$ $G$ abelian}
\label{subsec:gz-cyclic}

If $G/Z(G)$ is cyclic, then $G$ is abelian (and $G/Z(G)$ is trivial).

\subsection{N/C theorem}
\label{subsec:nc-theorem}

$H\le G$ $\Rightarrow$ $N_G(H)/C_G(H)\hookrightarrow\operatorname{Aut}(H)$. Special case: $G/Z(G)\cong\operatorname{Inn}(G)\le\operatorname{Aut}(G)$.

\subsection{Normalizer of a cyclic subgroup}
\label{subsec:normalizer-cyclic}

$\operatorname{ord}(g)=n$ $\Rightarrow$ $N_G(\langle g\rangle)/C_G(g)\hookrightarrow(\mathbb Z/n)^\times$. Consequence: $[N_G(\langle g\rangle):C_G(g)]\mid\varphi(n)$.

\subsection{Class size in a $p$-group}
\label{subsec:class-size-pgroup}

$|G|=p^n$, $g\notin Z(G)$ $\Rightarrow$ $|\operatorname{cl}(g)|=p^k$, $k\ge1$.

\medskip

\section{$S_n$ and $A_n$: sign, structure, conjugacy}
\label{sec:sn-an}

\subsection{Sign as a homomorphism}
\label{subsec:sign-homomorphism}

The unique $\operatorname{sgn}:S_n\to\{\pm1\}$ with $\operatorname{sgn}(\text{transposition})=-1$. $\ker\operatorname{sgn}=A_n$, $[S_n:A_n]=2$.

\subsection{Cycle type is a complete invariant for conjugacy in $S_n$}
\label{subsec:cycle-type}

$\sigma,\tau\in S_n$ are conjugate $\Leftrightarrow$ they have the same cycle type. Class size with type $1^{m_1}2^{m_2}\cdots n^{m_n}$: $\dfrac{n!}{\prod_k m_k!\,k^{m_k}}$.
$|C_{S_n}(\sigma)|=\prod_k m_k!\,k^{m_k}$.

\subsection{Conjugation formula}
\label{subsec:conjugation-formula}

$\rho\,(a_1\dots a_k)\,\rho^{-1}=(\rho(a_1)\dots\rho(a_k))$.

\subsection{Sign via cycle structure}
\label{subsec:sign-cycle}

$\operatorname{sgn}(\sigma)=(-1)^{n-c(\sigma)}=(-1)^{\sum(\lambda_i-1)}$, where $c(\sigma)$ is the number of cycles. A $k$-cycle has sign $(-1)^{k-1}$.

\subsection{Generators of $S_n$}
\label{subsec:generators-sn}

$S_n$ is generated by: all transpositions; adjacent transpositions; the pair $(1\,2)$ and $(1\,2\,\dots\,n)$.

\subsection{Generators of $A_n$}
\label{subsec:generators-an}

$A_n$ for $n\ge3$ is generated by $3$-cycles $(1\,2\,k)$, $k=3,\dots,n$.

\subsection{Splitting of $S_n$-classes in $A_n$}
\label{subsec:splitting-classes}

A class of $\sigma\in A_n$ in $S_n$ splits into two classes in $A_n$ $\Leftrightarrow$ the cycle type of $\sigma$ consists of pairwise distinct odd lengths.

\subsection{Simplicity of $A_n$ ($n\ge5$)}
\label{subsec:simplicity-an}

A group $G\neq\{e\}$ is \textit{simple} if it has no non-trivial normal subgroups. $A_n$ is simple for $n\ge5$.

\subsection{Normal subgroups of $S_n$}
\label{subsec:normal-sn}

$n=1,2$: $\{e\},S_n$. $n=3$: $\{e\},A_3,S_3$. $n=4$: $\{e\},V_4,A_4,S_4$. $n\ge5$: $\{e\},A_n,S_n$.

\subsection{Commutator}
\label{subsec:commutator-sn}

The \textit{commutator subgroup} $[G,G]=\langle g^{-1}h^{-1}gh:g,h\in G\rangle$ is the smallest normal subgroup with abelian quotient. $[S_n,S_n]=A_n$ ($n\ge2$). $[A_n,A_n]=A_n$ for $n\ge5$ (perfect). $[A_4,A_4]=V_4$.

\subsection{Solvability of $S_n$, $A_n$}
\label{subsec:solvability-sn}

A group $G$ is \textit{solvable} if its derived series $G^{(0)}=G$, $G^{(k+1)}=[G^{(k)},G^{(k)}]$ reaches $\{e\}$ finitely. $S_n$ and $A_n$ are solvable $\Leftrightarrow$ $n\le4$.

\medskip

\section{Sylow theorems and $p$-groups}
\label{sec:sylow}

\subsection{Sylow theorems}
\label{subsec:sylow}

$|G|=p^a m$, $\gcd(p,m)=1$.
\begin{itemize}
\item \textit{I (existence).} For each $0\le k\le a$ there is a subgroup of order $p^k$. A subgroup of order $p^a$ is a Sylow $p$-subgroup.
\item \textit{II (conjugacy).} All Sylow $p$-subgroups are conjugate. Any $p$-subgroup lies in some Sylow $p$-subgroup.
\item \textit{III (counting).} $n_p\equiv1\pmod p$, $n_p\mid m$, $n_p=[G:N_G(P)]$.
\end{itemize}

\subsection{Frattini argument}
\label{subsec:frattini}

$N\trianglelefteq G$, $P\in\operatorname{Syl}_p(N)$ $\Rightarrow$ $G=N\cdot N_G(P)$.

\subsection{$p$-groups are nilpotent}
\label{subsec:pgroup-nilpotent}

A group $G$ is \textit{nilpotent} if its upper central series $\{e\}=Z_0\le Z_1\le\cdots\le Z_k=G$ reaches $G$ in finitely many steps ($Z_{i+1}/Z_i=Z(G/Z_i)$). For $|G|=p^n$ this holds in at most $n$ steps. Consequence: in a $p$-group every proper $H$ is strictly smaller than $N_G(H)$.

\subsection{Maximal subgroup of a $p$-group}
\label{subsec:pgroup-maximal}

$|G|=p^n$, $M$ is maximal $\Rightarrow$ $M\trianglelefteq G$ and $[G:M]=p$.

\subsection{Frattini subgroup}
\label{subsec:frattini-subgroup}

$\Phi(G)=\bigcap M$ (intersection of all maximal subgroups). For a $p$-group: $\Phi(G)=[G,G]\cdot G^p$.
$G/\Phi(G)\cong(\mathbb F_p)^d$, $d$ is the minimal number of generators.

\subsection{Chain of normal subgroups in a $p$-group}
\label{subsec:pgroup-chain}

$|G|=p^n$ $\Rightarrow$ $\exists$ $\{e\}=N_0\lneq\cdots\lneq N_n=G$ with $|N_k|=p^k$ and $N_k\trianglelefteq G$.

\medskip

\section{Finite abelian groups}
\label{sec:finite-abelian}

\subsection{Structure theorem}
\label{subsec:structure-abelian}

Any finite abelian $G$:
\begin{itemize}
\item \textit{Invariant factors:} $G\cong\mathbb Z/d_1\oplus\cdots\oplus\mathbb Z/d_k$, $1<d_1\mid d_2\mid\cdots\mid d_k$.
\item \textit{Elementary divisors:} $G\cong\bigoplus\mathbb Z/p_i^{a_i}$ --- uniquely determined up to permutation.
\end{itemize}

\subsection{$p$-Sylow decomposition}
\label{subsec:p-sylow-abelian}

$G$ abelian $\Rightarrow$ $G=\bigoplus_{p\mid|G|}G_p$, where $G_p$ is the $p$-Sylow subgroup.

\subsection{$G[n]$ in an abelian group}
\label{subsec:g-n-abelian}

$G\cong\bigoplus\mathbb Z/p_i^{a_i}$ $\Rightarrow$ $|G[n]|=\prod_i\gcd(n,p_i^{a_i})$.

\subsection{Converse of Lagrange for abelian groups}
\label{subsec:converse-lagrange}

$G$ abelian, $d\mid|G|$ $\Rightarrow$ $\exists$ $H\le G$ of order $d$.

\subsection{Duality of characters}
\label{subsec:character-duality}

$G$ abelian. A \textit{character} is a homomorphism $\chi:G\to\mathbb C^\times$. $\widehat G=\operatorname{Hom}(G,\mathbb C^\times)\cong G$.
Orthogonality: $\sum_g\chi(g)\overline{\psi(g)}=|G|\cdot[\chi=\psi]$; $\sum_\chi\chi(g)\overline{\chi(h)}=|G|\cdot[g=h]$.

\subsection{Smith normal form}
\label{subsec:snf-abelian}

$M\in\mathbb Z^{m\times n}$ can be reduced to $\operatorname{diag}(d_1,\dots,d_r,0,\dots)$ with $d_1\mid d_2\mid\cdots\mid d_r$ by multiplication by $GL_m(\mathbb Z),GL_n(\mathbb Z)$. Application: if $G\cong\mathbb Z^n/M\mathbb Z^m$, then $G\cong\mathbb Z^{n-r}\oplus\bigoplus\mathbb Z/d_i$ (see also \ref{subsec:snf}).

\medskip

\section{Basic methods and techniques}
\label{sec:group-methods}

\begin{itemize}
\item \textit{Sylow counting.} $n_p\equiv1\pmod p$, $n_p\mid m$ often determines $n_p$. $n_p=1$ $\Rightarrow$ the Sylow subgroup is normal $\Rightarrow$ the group is not simple.
\item \textit{Element counting.} If $n_p$ is large, elements of order $p$ occupy much space, leaving little for others --- often $n_q=1$.
\item \textit{Class equation.} $|G|=|Z(G)|+\sum[G:C_G(g)]$ is the foundation of $p$-group induction.
\item \textit{Double counting.} The kernel of Burnside's lemma; counting pairs $(g,x)$ with $gx=x$.
\item \textit{Action on cosets.} Turns a subgroup $H$ into a homomorphism $G\to S_{[G:H]}$. The kernel is $\operatorname{core}_G(H)$, a normal subgroup inside $H$.
\item \textit{N/C theorem.} $N_G(H)/C_G(H)\hookrightarrow\operatorname{Aut}(H)$. Restricts the normalizer via automorphisms.
\item \textit{Reduction to $p$-parts.} In the abelian case the problem splits into independent $p$-components.
\end{itemize}
